% GENERATED FILE. Edit canonical lessons, then run npm run generate:books.
% canonical-source-hash: fnv1a64:895b20a37d131793
% canonical-lessons: JA-C66-kago, JA-C66-koma, JA-C66-ikada, JA-C66-fue, JA-C66-taiko

\chapter{Basket, Spinning top, Raft, Flute, Drum}
\label{ch:ja-basket-spinning-top-raft-flute-drum}
\hlchaptermodality{66}

\begin{chapteropening}
\textbf{By the end of this chapter:} \emph{I can say a basket, a spinning top, a raft, a flute and a drum.}

Complete the last lesson of chapter 66: I can say a basket, a spinning top, a raft, a flute and a drum.
\end{chapteropening}

\section[\texorpdfstring{\ja{かご} (kago)}{かご (kago)}]{\ja{かご} (kago) — a basket}

\label{lesson:JA-C66-kago}



\begin{quote}
Before the new one: say the Japanese for a potato, then the Japanese for burdock root.
\end{quote}

\subsection*{You'll want to know: \ja{かご}}
\textbf{\ja{かご}} — \emph{kago} — \textquotedblleft{}a basket\textquotedblright{}.

\begin{cousinweb}
\textbf{\ja{かご}} (\emph{kago}) is a native Japanese word, inherited rather than borrowed from Chinese, and its older origin is not settled. It is usually written in kana; its kanji, when used, is a kun'yomi attached later (JMdict, KANJIDIC2).
\end{cousinweb}

\begin{grammarlens}[title={What you've built}]
One more word for this chapter.
\end{grammarlens}

\subsection*{Your turn}
\begin{itemize}
\raggedright
  \item \emph{Say it:} \emph{kago}
  \item \emph{Say it:} \emph{kago}, once more
  \item \emph{Say it:} say \emph{kago}
  \item \emph{Recall:} say the Japanese for a potato, then the Japanese for burdock root, then say \emph{kago} again
\end{itemize}

\begin{tcolorbox}[breakable,colback=teal!4,colframe=teal!35!black,title={Before you move on}]
What does \ja{かご} mean? (\textquotedblleft{}A basket\textquotedblright{}.) Say it once more.
\end{tcolorbox}
\section[\texorpdfstring{\ja{こま} (koma)}{こま (koma)}]{\ja{こま} (koma) — a spinning top}

\label{lesson:JA-C66-koma}



\begin{quote}
Before the new one: say the Japanese for burdock root, then the Japanese for a basket.
\end{quote}

\subsection*{You'll want to know: \ja{こま}}
\textbf{\ja{こま}} — \emph{koma} — \textquotedblleft{}a spinning top\textquotedblright{}.

\begin{cousinweb}
\textbf{\ja{こま}} (\emph{koma}) is a native word, inherited rather than borrowed. JMdict writes it with characters for \textquotedblleft{}alone\textquotedblright{} and \textquotedblleft{}enjoy\textquotedblright{}, read as a whole, and marks it as usually written in kana. Its older origin is not settled.
\end{cousinweb}

\begin{grammarlens}[title={What you've built}]
One more word for this chapter.
\end{grammarlens}

\subsection*{Your turn}
\begin{itemize}
\raggedright
  \item \emph{Say it:} \emph{koma}
  \item \emph{Say it:} \emph{koma}, once more
  \item \emph{Say it:} say \emph{koma}
  \item \emph{Recall:} say the Japanese for burdock root, then the Japanese for a basket, then say \emph{koma} again
\end{itemize}

\begin{tcolorbox}[breakable,colback=teal!4,colframe=teal!35!black,title={Before you move on}]
What does \ja{こま} mean? (\textquotedblleft{}A spinning top\textquotedblright{}.) Say it once more.
\end{tcolorbox}
\section[\texorpdfstring{\ja{いかだ} (ikada)}{いかだ (ikada)}]{\ja{いかだ} (ikada) — a raft}

\label{lesson:JA-C66-ikada}



\begin{quote}
Before the new one: say the Japanese for a basket, then the Japanese for a spinning top.
\end{quote}

\subsection*{You'll want to know: \ja{いかだ}}
\textbf{\ja{いかだ}} — \emph{ikada} — \textquotedblleft{}a raft\textquotedblright{}.

\begin{cousinweb}
\textbf{\ja{いかだ}} (\emph{ikada}) is a native Japanese word, inherited rather than borrowed from Chinese, and its older origin is not settled. It is usually written in kana; its kanji, when used, is a kun'yomi attached later (JMdict, KANJIDIC2).
\end{cousinweb}

\begin{grammarlens}[title={What you've built}]
One more word for this chapter.
\end{grammarlens}

\subsection*{Your turn}
\begin{itemize}
\raggedright
  \item \emph{Say it:} \emph{ikada}
  \item \emph{Say it:} \emph{ikada}, once more
  \item \emph{Say it:} say \emph{ikada}
  \item \emph{Recall:} say the Japanese for a basket, then the Japanese for a spinning top, then say \emph{ikada} again
\end{itemize}

\begin{tcolorbox}[breakable,colback=teal!4,colframe=teal!35!black,title={Before you move on}]
What does \ja{いかだ} mean? (\textquotedblleft{}A raft\textquotedblright{}.) Say it once more.
\end{tcolorbox}
\section[\texorpdfstring{\ja{ふえ} (fue)}{ふえ (fue)}]{\ja{ふえ} (fue) — a flute}

\label{lesson:JA-C66-fue}



\begin{quote}
Before the new one: say the Japanese for a spinning top, then the Japanese for a raft.
\end{quote}

\subsection*{You'll want to know: \ja{ふえ}}
\textbf{\ja{ふえ}} — \emph{fue} — \textquotedblleft{}a flute\textquotedblright{}.

\begin{cousinweb}
\textbf{\ja{ふえ}} (\emph{fue}) is a native Japanese word, inherited rather than borrowed from Chinese, and its older origin is not settled. Its kanji was attached to it later as a kun'yomi, a native reading given to a Chinese character (JMdict, KANJIDIC2).
\end{cousinweb}

\begin{grammarlens}[title={What you've built}]
One more word for this chapter.
\end{grammarlens}

\subsection*{Your turn}
\begin{itemize}
\raggedright
  \item \emph{Say it:} \emph{fue}
  \item \emph{Say it:} \emph{fue}, once more
  \item \emph{Say it:} say \emph{fue}
  \item \emph{Recall:} say the Japanese for a spinning top, then the Japanese for a raft, then say \emph{fue} again
\end{itemize}

\begin{tcolorbox}[breakable,colback=teal!4,colframe=teal!35!black,title={Before you move on}]
What does \ja{ふえ} mean? (\textquotedblleft{}A flute\textquotedblright{}.) Say it once more.
\end{tcolorbox}
\section[\texorpdfstring{\ja{たいこ} (taiko)}{たいこ (taiko)}]{\ja{たいこ} (taiko) — a drum}

\label{lesson:JA-C66-taiko}



\begin{quote}
Before the new one: say the Japanese for a raft, then the Japanese for a flute.
\end{quote}

\subsection*{You'll want to know: \ja{たいこ}}
\textbf{\ja{たいこ}} — \emph{taiko} — \textquotedblleft{}a drum\textquotedblright{}.

\begin{cousinweb}
\textbf{\ja{たいこ}} (\emph{taiko}) is Sino-Japanese (kango): two Chinese characters read with their on'yomi, readings borrowed from Chinese. They are \emph{tai}, \textquotedblleft{}big, great\textquotedblright{} (Middle Chinese \emph{thajH}) and \emph{ko}, \textquotedblleft{}drum\textquotedblright{} (\emph{kuX}) (Unihan, KANJIDIC2; Middle Chinese in Baxter's transcription).
\end{cousinweb}

\begin{grammarlens}[title={What you've built}]
One more word for this chapter.
\end{grammarlens}

\subsection*{Your turn}
\begin{itemize}
\raggedright
  \item \emph{Say it:} \emph{taiko}
  \item \emph{Say it:} \emph{taiko}, once more
  \item \emph{Say it:} say \emph{taiko}
  \item \emph{Recall:} say the Japanese for a raft, then the Japanese for a flute, then say \emph{taiko} again
\end{itemize}

\begin{tcolorbox}[breakable,colback=teal!4,colframe=teal!35!black,title={Before you move on}]
What does \ja{たいこ} mean? (\textquotedblleft{}A drum\textquotedblright{}.) Say it once more.
\end{tcolorbox}
